% Adapted from the author's public Erdos Paper IV v1.23 series_template.tex.
% Layout provenance is recorded in 02_validation; the Markdown is semantic authority.
\documentclass[english,11pt,a4paper]{article}
\usepackage[margin=1in]{geometry}
\usepackage{fontspec}
\setmainfont{lmroman10-regular.otf}[BoldFont=lmroman10-bold.otf,ItalicFont=lmroman10-italic.otf,BoldItalicFont=lmroman10-bolditalic.otf]
\setsansfont{lmsans10-regular.otf}[BoldFont=lmsans10-bold.otf]
\setmonofont{lmmono10-regular.otf}[Scale=MatchLowercase]
\usepackage{amsmath,amssymb}
\providecommand{\mathscr}{\mathcal}
\usepackage[english]{babel}
\usepackage{microtype,parskip}
\usepackage{longtable,booktabs,array,calc}
\usepackage{xcolor,xurl}
\usepackage{hyperref,bookmark}
\usepackage{etoolbox,needspace}
\setcounter{secnumdepth}{-\maxdimen}
\setlength{\emergencystretch}{3em}
\providecommand{\tightlist}{\setlength{\itemsep}{0pt}\setlength{\parskip}{0pt}}
\newcounter{none}
\AtBeginEnvironment{longtable}{\Needspace{8\baselineskip}\small}
\pretocmd{\section}{\Needspace{5\baselineskip}}{}{}
\pretocmd{\subsection}{\Needspace{7\baselineskip}}{}{}
\pretocmd{\subsubsection}{\Needspace{6\baselineskip}}{}{}
\urlstyle{same}

\usepackage{graphicx}
\providecommand{\pandocbounded}[1]{#1}

\usepackage{float}

\newcommand{\KL}{\operatorname{KL}}
\clubpenalty=10000\widowpenalty=10000\displaywidowpenalty=10000
\hypersetup{pdftitle={Formal companion: entropy and large deviations for Latin squares},pdfauthor={Juan Pablo Traverso Gianini},pdflang={en},colorlinks=true,linkcolor=blue,urlcolor=blue,citecolor=blue}
\title{Formal companion: entropy and large deviations for Latin squares}
\author{Juan Pablo Traverso Gianini\\\small Independent researcher, Santiago, Chile\\\small \href{mailto:jtraverso@gmail.com}{jtraverso@gmail.com}\\\small \href{https://orcid.org/0009-0003-6068-4096}{ORCID: 0009-0003-6068-4096}}
\date{}\begin{document}\raggedbottom\maketitle
\textbf{Preprint 1.0 --- 16 November 2026.}

\textbf{Repository:} \url{https://github.com/jtraverso/latinon-entropy-ldp}.

\subsection{1. Scope and mathematical interface}\label{scope-and-mathematical-interface}

This companion records how the statements of the manuscript are represented in Lean. It is a guide to the delivered sources and evidence, not a substitute for the written proofs. The model uses three uniform two-coordinate marginals and retains the measure-preserving order map. Its distance has two order terms and a distribution term testing symbol intervals. Relative entropy takes values in the extended nonnegative reals. Converting it to a real number is justified only after a finite-rate hypothesis.

The formal statements are the protected contracts in the audit record. The upper deviation endpoint uses \texttt{ENNReal.ofReal\ t}; this is the nonstrict floor max(t,0) in the manuscript. The lower endpoint requires a nonnegative finite witness cost. Neither endpoint presupposes a law on an auxiliary collection: both concern the pushforward of the actual uniform finite Latin-square law.

\subsection{2. Principal statement correspondence}\label{principal-statement-correspondence}

\Needspace{7\baselineskip}\textbf{Theorem 2.1 / Teorema 2.1}

\path{LatinonEntropy.closed_uniform_latinon_upper}

Source: \path{lean/LatinonBridge/LatinonGlobalDeviations.lean}, line 20.

\path{LatinonEntropy.open_uniform_latinon_lower}

Source: \path{lean/LatinonBridge/LatinonGlobalDeviations.lean}, line 46.

\Needspace{7\baselineskip}\textbf{Theorem 2.2 / Teorema 2.2}

\path{LatinonEntropy.finite_rate_standard_pattern_log_count}

Source: \path{lean/LatinonPatternConjecture.lean}, line 12.

\path{LatinonEntropy.empty_standard_pattern_fiber_eventually_zero}

Source: \path{lean/LatinonPatternConjecture.lean}, line 42.

\Needspace{7\baselineskip}\textbf{Theorem 3.1 / Teorema 3.1}

\path{LatinonEntropy.finite_latin_tensor_lower_bound}

Source: \path{lean/LatinTensorLowerBound.lean}, line 51.

\Needspace{7\baselineskip}\textbf{Lemma 4.1 / Lema 4.1}

\path{LatinonEntropy.correction_lemma}

Source: \path{lean/QuasirandomCorrection.lean}, line 143.

Proposition 2.3 is the written identification with the published GHHS representations and cut-distance definition. Its disintegration and change-of-variables proof is given in the manuscript; it is not advertised as a new Lean declaration. The formal source uses the measure representation, inverse coordinate pushforwards and two order terms. The manuscript notation does not rename any Lean declaration.

\textbf{Good rate in Theorem 2.1}

\path{LatinonEntropy.latinon_rate_lowerSemicontinuous}

Source: \path{lean/LatinonFullDyadicKL.lean}, line 42.

\path{LatinonEntropy.latinon_full_cut_compact}

Source: \path{lean/LatinonBridge/LatinonCompactness.lean}, line 60.

\path{LatinonEntropy.latinon_rate_sublevels_compact}

Source: \path{lean/LatinonBridge/LatinonGlobalDeviations.lean}, line 15.

These three properties are queried together with the eight original endpoints and the strengthened correction theorem in \path{ReviewAxioms.lean}. All twelve current queries report only standard foundational axioms. The original eleven outputs are also retained with the clean reconstruction in \path{preprints/LATINON/02_validation/CLEAN_REBUILD_AUDIT.json}.

The finite tensor theorem quantifies the order threshold before the profile. The correction theorem has no residual-diameter or potential certificate among its hypotheses. Its explicit box constants and linear coefficient 3 are preserved in Lemma 4.1.

The diameter-six argument is also present internally as \path{LatinonEntropy.ScaledResidualData.diameter_six} in \path{lean/ResidualDiameter.lean}. Its absence from the hypotheses of \path{correction_lemma} means it is derived internally, not that it exists only on paper. Review Lemma 4.1 as a written proof separately from checking its correspondence with that source argument.

\textbf{Refinement of Lemma 4.1.}

\path{LatinonEntropy.ScaledResidualData.diameter_five}

\path{LatinonEntropy.correction_lemma_five}

Both declarations are in \path{lean/CorrectionFive.lean}. The transpose reverses residual arcs; its four-step row paths give four-step column paths in the original graph. Opposite-side paths use at most five edges, giving the potential radius and correction coefficient \(5/2\). These declarations strengthen the original coefficient-three result without changing its consumers or the numbered theorem correspondence.

\subsection{3. Dependency ledger}\label{dependency-ledger}

The core ledger has 23 steps. The following list gives the declaration and dependencies at each step. The current compiled-source correspondence is in \path{publication/PROOF_LEDGER.json}. Original source-bound certificates retained inside the Lean archive are historical evidence; the completed reconstruction report certifies the delivered selection.

\Needspace{6\baselineskip}\textbf{1. Model and full cut distance.}

\path{LatinonEntropy.LatinonMeasure}

Ledger dependencies: none.

\Needspace{6\baselineskip}\textbf{2. Deterministic correction.}

\path{LatinonEntropy.correction_lemma}

Ledger dependencies: none.

\Needspace{6\baselineskip}\textbf{3. Permanents and marginals.}

\path{LatinonEntropy.permanental_marginal_l1_le}

Ledger dependencies: none.

\Needspace{6\baselineskip}\textbf{4. Latin process and stopping.}

\path{LatinonEntropy.profileProcessPolicy}

Ledger dependencies: 2, 3.

\Needspace{6\baselineskip}\textbf{5. Defect recurrence.}

\path{LatinonEntropy.profile_expected_defect_recurrence}

Ledger dependencies: 4.

\Needspace{6\baselineskip}\textbf{6. Gronwall stability.}

\path{LatinonEntropy.profile_expected_defect_uniform}

Ledger dependencies: 5.

\Needspace{6\baselineskip}\textbf{7. Uniform survival.}

\path{LatinonEntropy.eventually_uniform_profile_survival}

Ledger dependencies: 6.

\Needspace{6\baselineskip}\textbf{8. Uniform cut closeness.}

\path{LatinonEntropy.eventually_uniform_profile_closeness}

Ledger dependencies: 7.

\Needspace{6\baselineskip}\textbf{9. Entropy of the full process.}

\path{LatinonEntropy.profile_history_entropy_lower}

Ledger dependencies: 6, 7.

\Needspace{6\baselineskip}\textbf{10. Tensor lower counting.}

\path{LatinonEntropy.finite_latin_tensor_lower_bound}

Ledger dependencies: 8, 9.

\Needspace{6\baselineskip}\textbf{11. Tensor-Latinon bridge.}

\path{LatinonEntropy.finite_latin_grid_metric_ball_lower_bound}

Ledger dependencies: 1, 10.

\Needspace{6\baselineskip}\textbf{12. Partition upper bound.}

\path{LatinonEntropy.exists_latinon_partition_ball_count_upper}

Ledger dependencies: 1.

\Needspace{6\baselineskip}\textbf{13. Extended KL rate.}

\path{LatinonEntropy.latinon_rate_eq_iSup_dyadic}

Ledger dependencies: 12.

\Needspace{6\baselineskip}\textbf{14. Fixed center and order.}

\path{LatinonEntropy.exists_bounded_latinon_profile_approximation}

Ledger dependencies: 11, 13.

\Needspace{6\baselineskip}\textbf{15. Regularization and local entropy.}

\path{LatinonEntropy.finite_rate_latinon_entropy_count}

Ledger dependencies: 14.

\Needspace{6\baselineskip}\textbf{16. Full-distance compactness.}

\path{LatinonEntropy.latinon_full_cut_compact}

Ledger dependencies: 1.

\Needspace{6\baselineskip}\textbf{17. Global LDP.}

\path{LatinonEntropy.closed_uniform_latinon_upper}

Ledger dependencies: 13, 15, 16.

\Needspace{6\baselineskip}\textbf{18. Contraction and observables.}

\path{LatinonEntropy.latinonContractedRate_attained}

Ledger dependencies: 17.

\Needspace{6\baselineskip}\textbf{19. Densities, continuity and normalization.}

\path{LatinonEntropy.sum_latinonPatternDensity}

Ledger dependencies: 1.

\Needspace{6\baselineskip}\textbf{20. Finite coordinate bridge.}

\path{LatinonEntropy.finite_matrix_pattern_embedding_error}

Ledger dependencies: 19.

\Needspace{6\baselineskip}\textbf{21. Binomial normalization.}

\path{LatinonEntropy.standard_pattern_profile_uniform_error}

Ledger dependencies: 20.

\Needspace{6\baselineskip}\textbf{22. Counting under closed constraints.}

\path{LatinonEntropy.finite_rate_standard_pattern_entropy_count}

Ledger dependencies: 18, 21.

\Needspace{6\baselineskip}\textbf{23. Logarithmic count and empty fiber.}

\path{LatinonEntropy.finite_rate_standard_pattern_log_count}

Ledger dependencies: 22, 16.

\subsection{4. Applications and source assemblies}\label{applications-and-source-assemblies}

The application assembly is \path{LatinonSeriesApplications.lean}. It contains the finite CDR definitions and the Wilson pair-footprint argument, then the factorization correspondence, packing replacement, local entropy products and exact type preservation. Its eight original inputs are retained under \path{lean/evidence/applications}. The core pattern assembly similarly retains its 18 proof inputs under \path{lean/evidence/proof-inputs}. An input being retained does not assert that it is compiled as an independent module: the assembled modules are the compilation targets.

\textbf{A1.} \path{finite_rate_factorization_entropy_count}

\textbf{A2.} \path{finite_rate_factorization_pattern_log_count}

\textbf{A3.} \path{replaceSteinerPacking_bicoloured}

\textbf{A4.} \path{wilsonLatinBlockTemplate}

\textbf{A5.} \path{card_restrictedRealizations}

\textbf{A6.} \path{wilson_distinguished_latinon_blocks_entropy_count}

\textbf{A7.} \path{wilson_completed_exact_type_count_eq_base}

\textbf{A8.} \path{residual_subset_completed}

The CDR endpoints quantify an already completed base system containing the distinguished packing. They do not discharge the existence of that base from class-size inequalities alone. The final families are actual block sets, not marked copies. The local squares are recovered from fixed row-column pairs, so their product count has no multiplicity loss. Block-type counts refer to color multisets and are equal to the base counts; fixed configurations in the residue are retained, without asserting that the total number of such configurations is unchanged.

\subsection{5. Toolchain, imports and reconstruction}\label{toolchain-imports-and-reconstruction}

The current project pins Lean 4.35.0, Mathlib commit \path{56a005855ea17e0b0958f73ec801ae140b0a1679} and eight further ecosystem repositories. It declares 310 local module roots. The canonical entry \path{LatinonFormalization} imports the applications assembly and \path{CorrectionFive}; its local import closure has 299 modules. The source manifest lists all modules and selected file hashes.

Extract \path{preprints/LATINON/05_formalization/Latinon_Lean_v1.0.zip}, which contains the project \path{lean/} and upstream licenses. Install the pinned toolchain and external packages, then run \texttt{lake\ build\ LatinonFormalization} and \texttt{lake\ env\ lean\ ReviewAxioms.lean} sequentially from that project. The library build covers 310 configured targets. The original 309-target reconstruction began without project objects; the subsequent incremental build reused those objects and pinned external caches.

Both builds completed on 16 November 2026. The current twelve-query file and its successful output are delivered with the incremental evidence. All reported axiom sets are contained in \path{propext}, \path{Classical.choice} and \path{Quot.sound}. The archive and logs are bound by SHA-256; read-only package verification does not invoke Lean.

Run \texttt{python\ scripts/review\_axioms.py} from the publication repository to verify all twelve recorded queries in one command. This checks certified log and source bindings. For a fresh Lean execution, use the query command above.

\subsection{6. Provenance and review limits}\label{provenance-and-review-limits}

The original clean snapshot is identified by \path{7b603942c8d078db6c0a5c8729df7012b7786db8} and is retained separately. The current source freeze adds one module and changes the canonical entry, query file and library configuration. Its source manifest is regenerated for the delivered ZIP. \path{preprints/LATINON/02_validation/v1.0/FREEZE.json} binds the current archive and entry hashes; \path{BASELINE_FREEZE.json} binds the original evidence. From the repository root, \texttt{python\ scripts/verify\_repository.py} checks both without private Git objects.

Nine graphon modules originate from Cameron Freer's graphon project at commit \path{9f7be59fa754d260a544b4cfd83d6a5b94f7552e}. The retained provenance file lists removed blueprint metadata and compatibility changes. The finite CDR source cut originates from \path{jtraverso/cdr-palomar-kernel} at commit \path{5d1e6fe301e5fc01c92f6145aa34fddfe9c2bd10}. The basic definitions are byte-identical; Wilson compatibility edits and the Latin-array representation change are enumerated. This selection does not import the former absorption existence library.

The internal audit follows the CDR Paper III protocol at series commit \path{d12d980f8c86d8795e9faa0c6bb274ee936c7c8e}. It separates reconstruction, axiom checks, protected statements, targeted written-proof review, bilingual parity, rendered documents and archive integrity. The edition audit is \path{publication/LATINON_RELEASE_AUDIT.json}. It is an internal preparation check, not independent peer review or an exhaustive priority certification.

\end{document}
